% ============================================================
%  Chapter 2 — Robot Bodies: Arms, Wheels, and Legs
%  Robot World: Meet the Machines That Help Us
% ============================================================

\chapter{Robot Bodies \textemdash{} Arms, Wheels, and Legs}

\chapteropening{%
  Humans have arms, legs, and hands. Robots have joints, wheels, and grippers.
  Let us explore how robot bodies are built to move.
}
\chapterbadge{Meet Bolt \textemdash{} the factory arm!}

% --- Introduction ---
\section*{Built for a Job}

A robot's body is designed for its task. A factory robot arm does not need
legs \textemdash{} it stays in one place and reaches. A delivery robot does not
need arms \textemdash{} it just needs wheels and a cargo box.

\character{Bolt}{My body has six joints, just like your shoulder, elbow,
  and wrist \textemdash{} but doubled! That means I can twist and reach in ways
  no human arm can.}

\begin{picturethis}
  Think of a crane at a building site. It is tall and has a long arm that
  swings. A robot arm is like a tiny, super-precise crane that sits on a
  table and can move to the exact same spot a thousand times without missing.
\end{picturethis}

\sectionrule

% --- Moving parts ---
\section*{How Robots Move}

Every moving robot uses some combination of these:

\begin{itemize}
  \item \term{Motors} \textemdash{} electric motors spin wheels or rotate joints.
  \item \term{Joints} \textemdash{} hinges that let parts bend or twist.
  \item \term{Gears} \textemdash{} toothed wheels that change speed and force.
  \item \term{Wheels} \textemdash{} great for flat surfaces; fast and efficient.
  \item \term{Legs} \textemdash{} better for rough terrain but harder to control.
  \item \term{Tracks} \textemdash{} like tank treads; good on mud, rubble, or sand.
  \item \term{Propellers} \textemdash{} spinning blades that push air down to make drones fly.
\end{itemize}

\begin{funfact}
  Boston Dynamics' robot dog ``Spot'' has 12 motors \textemdash{}
  three in each leg \textemdash{} and can climb stairs, open doors, and even dance!
\end{funfact}

\sectionrule

% --- Grippers ---
\section*{Robot Hands and Grippers}

Robots need to grab things too. But they do not all use fingers.

\begin{enumerate}
  \item \textbf{Pincer grippers} \textemdash{} two fingers that squeeze together (like tongs).
  \item \textbf{Suction cups} \textemdash{} use a vacuum to stick to flat surfaces (like glass).
  \item \textbf{Magnetic grippers} \textemdash{} attract metal objects.
  \item \textbf{Soft grippers} \textemdash{} inflate like balloons to wrap gently around odd shapes.
  \item \textbf{Multi-finger hands} \textemdash{} five fingers controlled independently, the closest to a human hand.
\end{enumerate}

\begin{bigidea}
  Engineers choose the simplest gripper that can do the job.
  More fingers means more complexity \textemdash{} and more things that can go wrong.
\end{bigidea}

\sectionrule

% --- Materials ---
\section*{What Are Robots Made Of?}

\begin{itemize}
  \item \textbf{Aluminium and steel} \textemdash{} strong but light for arms and frames.
  \item \textbf{Plastic and carbon fibre} \textemdash{} lightweight for drones.
  \item \textbf{Silicone and rubber} \textemdash{} flexible, used in soft robots.
  \item \textbf{Titanium} \textemdash{} extremely tough, used in space robots.
\end{itemize}

\character{Sprocket}{I am made mostly of aluminium. That keeps me light
  enough to roll over sand dunes without sinking!}

\sectionrule

% --- Experiment ---
\begin{experiment}
  \textbf{Build a simple gripper.}
  \begin{enumerate}
    \item Take two wooden lolly sticks and a rubber band.
    \item Place a small sponge between the tips.
    \item Wrap the rubber band around the other end to hold them together.
    \item Squeeze the sticks apart \textemdash{} the sponge tips close like a pincer.
    \item Try picking up a cotton ball, a marble, and a pencil. Which is hardest?
  \end{enumerate}
  \textit{This is how a simple pincer gripper works!}
\end{experiment}

\sectionrule

% --- Diagram ---
\begin{diagram}[Parts of a typical robot arm.]
  \begin{tikzpicture}[>=Stealth]
    % Base
    \draw[fill=MetalGray!30,draw=SteelBlue,line width=1pt,rounded corners=2pt]
      (-1.5,-0.3) rectangle (1.5,0.3);
    \node[font=\small\bfseries] at (0,0) {Base};
    % Arm segment 1
    \draw[fill=RoboGreen!20,draw=RoboGreen!80!black,line width=1pt,rounded corners=2pt]
      (-0.4,0.3) rectangle (0.4,2.3);
    \node[font=\small,rotate=90] at (0,1.3) {Segment 1};
    % Joint
    \fill[WarmOrange] (0,2.3) circle (0.25);
    \node[font=\tiny\bfseries,right] at (0.3,2.3) {Joint};
    % Arm segment 2
    \draw[fill=RoboGreen!20,draw=RoboGreen!80!black,line width=1pt,rounded corners=2pt]
      (-0.35,2.55) rectangle (0.35,4.2);
    \node[font=\small,rotate=90] at (0,3.4) {Segment 2};
    % Gripper
    \draw[fill=SunYellow!30,draw=SunYellow!70!black,line width=1pt]
      (-0.5,4.2) -- (-0.7,4.9) -- (-0.2,4.9) -- (-0.1,4.2) -- cycle;
    \draw[fill=SunYellow!30,draw=SunYellow!70!black,line width=1pt]
      (0.5,4.2) -- (0.7,4.9) -- (0.2,4.9) -- (0.1,4.2) -- cycle;
    \node[font=\small\bfseries,above] at (0,4.95) {Gripper};
  \end{tikzpicture}
\end{diagram}

\sectionrule


% --- Robot power ---
\section*{What Powers a Robot?}

Every motor, sensor, and controller needs \term{energy}. Most robots get
their energy from \term{batteries} --- portable packs of stored electricity.

\begin{itemize}
  \item \textbf{Rechargeable batteries} --- the most common; used in drones, vacuums, and rovers.
  \item \textbf{Mains power} --- factory robots often plug into the wall for unlimited energy.
  \item \textbf{Solar panels} --- some space robots collect energy from sunlight.
  \item \textbf{Charging docks} --- home robots return to a base station when running low.
\end{itemize}

\begin{picturethis}
  Imagine going to school without eating breakfast or lunch.
  By afternoon you would have no energy to run or think.
  A robot with a low battery is exactly the same --- it slows down
  and eventually stops. Good robot designers plan carefully
  so their robot always has enough energy to finish the job
  \emph{and} get back to its charger.
\end{picturethis}

\begin{funfact}
  A typical drone battery lasts only 20--30 minutes.
  That is why delivery drones must fly fast and land quickly!
\end{funfact}

\sectionrule

\begin{quickcheck}
  \begin{enumerate}
    \item What is the difference between a motor and a gear?
    \item Name two types of gripper and when you might use each one.
    \item Why do drones use propellers instead of wheels?
  \end{enumerate}
\end{quickcheck}

\sectionrule
% --- New Words ---
\begin{newwords}
  \begin{description}[labelwidth=3.8cm, leftmargin=4.0cm]
    \item[Motor] A device that converts electricity into spinning motion.
    \item[Joint] A connection between two parts that allows bending or rotating.
    \item[Gear] A toothed wheel that transfers force between shafts.
    \item[Gripper] A robot hand designed to grab objects.
    \item[Degrees of freedom] The number of independent ways a robot can move.
    \item[Chassis] The main frame or body of a robot.
  \end{description}
\end{newwords}
